\documentclass[../../main.tex]{subfiles}% 注意这里的文件路径不能用 ./main.tex ,否则用latexmk编译子文件会报错
\graphicspath{{\subfix{./image/}}{\subfix{../../image/}}} % 指定图片引用路径,后续可以直接使用图片文件名
% 如果图片存放在上级目录,则使用 ../../image/ 作为备用路径

% 例如:
% \begin{figure}[H]
% \centering
% \includegraphics[width=0.3\textwidth]{图.png}
% \caption{}
% \label{figure:图}
% \end{figure}
% 注意:上述\label{}一定要放在\caption{}之后,否则引用图片序号会只会显示??.

\begin{document}

\section{第二型曲面积分}

\begin{definition}
设连通曲面 $S$ 上到处都有连续变动的切平面或法线，$M$ 为曲面 $S$ 上的一点，曲面在 $M$ 处的法线有两个方向: 当取定其中一个指向为正方向时，则另一个指向就是负方向. 设 $M_0$ 为 $S$ 上任一点，$L$ 为 $S$ 上任一经过点 $M_0$，且不超出 $S$ 边界的闭曲线. 又设 $M$ 为动点，它在 $M_0$ 处与 $M_0$ 有相同的法线方向，且有如下特性: 当 $M$ 从 $M_0$ 出发沿 $L$ 连续移动，这时作为曲面上的点 $M$，它的法线方向也连续地变动. 最后当 $M$ 沿 $L$ 回到 $M_0$ 时，若这时 $M$ 的法线方向仍与 $M_0$ 的法线方向相一致，则说这曲面 $S$ 是双侧曲面；若与 $M_0$ 的法线方向相反，则说 $S$ 是\textbf{单侧曲面}.
\end{definition}

\begin{definition}
设 $P, Q, R$ 为定义在双侧曲面 $S$ 上的函数，在 $S$ 所指定的一侧作分割 $T$，它把 $S$ 分为 $n$ 个小曲面 $S_1, S_2, \cdots, S_n$，分割 $T$ 的细度 $\|T\| = \max\limits_{1\leqslant i\leqslant n} \{S_i \text{的直径}\}$，以 $\Delta S_{i_{yz}}, \Delta S_{i_{zx}}, \Delta S_{i_{xy}}$ 分别表示 $S_i$ 在三个坐标面上的投影区域的面积，它们的符号由 $S_i$ 的方向来确定. 若 $S_i$ 的法线正向与 $z$ 轴正向成锐角时，$S_i$ 在 $xy$ 平面的投影区域的面积 $\Delta S_{i_{xy}}$ 为正. 反之，若 $S_i$ 法线正向与 $z$ 轴正向成钝角时，它在 $xy$ 平面的投影区域的面积 $\Delta S_{i_{xy}}$ 为负. 在各个小曲面 $S_i$ 上任取一点 $(\xi_i, \eta_i, \zeta_i)$. 若
\begin{align*}
\lim\limits_{\|T\|\to 0} \sum\limits_{i=1}^n P(\xi_i, \eta_i, \zeta_i) \Delta S_{i_{yz}} + \lim\limits_{\|T\|\to 0} \sum\limits_{i=1}^n Q(\xi_i, \eta_i, \zeta_i) \Delta S_{i_{zx}} + \lim\limits_{\|T\|\to 0} \sum\limits_{i=1}^n R(\xi_i, \eta_i, \zeta_i) \Delta S_{i_{xy}}
\end{align*}
存在，且与曲面 $S$ 的分割 $T$ 和 $(\xi_i, \eta_i, \zeta_i)$ 在 $S_i$ 上的取法无关，则称此极限为函数 $P, Q, R$ 在曲面 $S$ 所指定的一侧上的\textbf{第二型曲面积分}，记作
\begin{align*}
\iint_S P(x,y,z) \mathrm{d}y\mathrm{d}z + Q(x,y,z) \mathrm{d}z\mathrm{d}x + R(x,y,z) \mathrm{d}x\mathrm{d}y
\end{align*}
或
\begin{align*}
\iint_S P(x,y,z) \mathrm{d}y\mathrm{d}z + \iint_S Q(x,y,z) \mathrm{d}z\mathrm{d}x + \iint_S R(x,y,z) \mathrm{d}x\mathrm{d}y.
\end{align*}
\end{definition}
\begin{remark}
据此定义，某流体以速度 $\bm{v} = (P, Q, R)$ 在单位时间内从曲面 $S$ 的负侧流向正侧的总流量
\begin{align*}
E = \iint_S P(x,y,z) \mathrm{d}y\mathrm{d}z + Q(x,y,z) \mathrm{d}z\mathrm{d}x + R(x,y,z) \mathrm{d}x\mathrm{d}y.
\end{align*}
又若，空间的磁场强度为 $(P(x,y,z), Q(x,y,z), R(x,y,z))$，则通过曲面 $S$ 的磁通量(磁力线总数)
\begin{align*}
H = \iint_S P(x,y,z) \mathrm{d}y\mathrm{d}z + Q(x,y,z) \mathrm{d}z\mathrm{d}x + R(x,y,z) \mathrm{d}x\mathrm{d}y.
\end{align*}
若以 $-S$ 表示曲面 $S$ 的另一侧，由定义易得
\begin{align*}
\iint_{-S} P \mathrm{d}y\mathrm{d}z + Q \mathrm{d}z\mathrm{d}x + R \mathrm{d}x\mathrm{d}y = -\iint_S P \mathrm{d}y\mathrm{d}z + Q \mathrm{d}z\mathrm{d}x + R \mathrm{d}x\mathrm{d}y.
\end{align*}
\end{remark}

\begin{proposition}
与第二型曲线积分一样，第二型曲面积分也有如下一些性质:
\begin{enumerate}[(1)]
\item 若 $\iint_S P_i \mathrm{d}y\mathrm{d}z + Q_i \mathrm{d}z\mathrm{d}x + R_i \mathrm{d}x\mathrm{d}y$ ($i=1, 2, \cdots, k$) 存在，则有
\begin{align*}
\iint_S \left(\sum\limits_{i=1}^k c_i P_i\right) \mathrm{d}y\mathrm{d}z + \left(\sum\limits_{i=1}^k c_i Q_i\right) \mathrm{d}z\mathrm{d}x + \left(\sum\limits_{i=1}^k c_i R_i\right) \mathrm{d}x\mathrm{d}y = \sum\limits_{i=1}^k c_i \iint_S P_i \mathrm{d}y\mathrm{d}z + Q_i \mathrm{d}z\mathrm{d}x + R_i \mathrm{d}x\mathrm{d}y,
\end{align*}
其中 $c_i$ ($i=1, 2, \cdots, k$) 是常数.
\item 若曲面 $S$ 是由两两无公共内点的曲面块 $S_1, S_2, \cdots, S_k$ 所组成，且
\begin{align*}
\iint_{S_i} P \mathrm{d}y\mathrm{d}z + Q \mathrm{d}z\mathrm{d}x + R \mathrm{d}x\mathrm{d}y \quad (i=1, 2, \cdots, k)
\end{align*}
存在，则有
\begin{align*}
\iint_S P \mathrm{d}y\mathrm{d}z + Q \mathrm{d}z\mathrm{d}x + R \mathrm{d}x\mathrm{d}y = \sum\limits_{i=1}^k \iint_{S_i} P \mathrm{d}y\mathrm{d}z + Q \mathrm{d}z\mathrm{d}x + R \mathrm{d}x\mathrm{d}y.
\end{align*}
\end{enumerate}
\end{proposition}

\begin{theorem}
设 $R$ 是定义在光滑曲面
\begin{align*}
S: z = z(x,y), \quad (x,y) \in D_{xy}
\end{align*}
上的连续函数，以 $S$ 的上侧为正侧(这时 $S$ 的法线方向与 $z$ 轴正向成锐角)，则有
\begin{align}
\iint_S R(x,y,z) \mathrm{d}x\mathrm{d}y = \iint_{D_{xy}} R(x,y,z(x,y)) \mathrm{d}x\mathrm{d}y. \label{eq:22-2}
\end{align}
\end{theorem}
\begin{proof}
由第二型曲面积分定义
\begin{align*}
\iint_S R(x,y,z) \mathrm{d}x\mathrm{d}y = \lim\limits_{\|T\|\to 0} \sum\limits_{i=1}^n R(\xi_i, \eta_i, \zeta_i) \Delta S_{i_{xy}} = \lim\limits_{d\to 0} \sum\limits_{i=1}^n R(\xi_i, \eta_i, z(\xi_i, \eta_i)) \Delta S_{i_{xy}},
\end{align*}
这里 $d = \max \{S_{i_{xy}} \text{的直径}\}$. 显然由 $\|T\| = \max \{S_i \text{的直径}\} \to 0$ 立刻可推得 $d \to 0$. 由于 $R$ 在 $S$ 上连续，$z$ 在 $D_{xy}$ 上连续(曲面光滑)，据复合函数的连续性，$R(x,y,z(x,y))$ 也是 $D_{xy}$ 上的连续函数. 由二重积分的定义
\begin{align*}
\iint_{D_{xy}} R(x,y,z(x,y)) \mathrm{d}x\mathrm{d}y = \lim\limits_{d\to 0} \sum\limits_{i=1}^n R(\xi_i, \eta_i, z(\xi_i, \eta_i)) \Delta S_{i_{xy}}.
\end{align*}
所以
\begin{align*}
\iint_S R(x,y,z) \mathrm{d}x\mathrm{d}y = \iint_{D_{xy}} R(x,y,z(x,y)) \mathrm{d}x\mathrm{d}y.
\end{align*}
\end{proof}

\begin{corollary}
类似地，当 $P$ 在光滑曲面
\begin{align*}
S: x = x(y,z), \quad (y,z) \in D_{yz}
\end{align*}
上连续时，有
\begin{align}
\iint_S P(x,y,z) \mathrm{d}y\mathrm{d}z = \iint_{D_{yz}} P(x(y,z), y, z) \mathrm{d}y\mathrm{d}z, \label{eq:22-3}
\end{align}
这里 $S$ 是以 $S$ 的法线方向与 $x$ 轴的正向成锐角的那一侧为正侧. 当 $Q$ 在光滑曲面
\begin{align*}
S: y = y(z,x), \quad (z,x) \in D_{zx}
\end{align*}
上连续时，有
\begin{align}
\iint_S Q(x,y,z) \mathrm{d}z\mathrm{d}x = \iint_{D_{zx}} Q(x, y(z,x), z) \mathrm{d}z\mathrm{d}x, \label{eq:22-4}
\end{align}
这里 $S$ 是以 $S$ 的法线方向与 $y$ 轴的正向成锐角的那一侧为正侧.
\end{corollary}

\begin{example}
计算 $\iint_S xyz \mathrm{d}x\mathrm{d}y$，其中 $S$ 是球面 $x^2 + y^2 + z^2 = 1$ 在 $x\geqslant 0, y\geqslant 0$ 部分并取球面外侧.
\end{example}
\begin{solution}
曲面 $S$ 在第一、五卦限部分的方程分别为
\begin{align*}
S_1: z_1 &= \sqrt{1 - x^2 - y^2}, \\
S_2: z_2 &= -\sqrt{1 - x^2 - y^2},
\end{align*}
它们在 $xy$ 平面上的投影区域都是单位圆在第一象限部分. 依题意，积分是沿 $S_1$ 的上侧和 $S_2$ 的下侧进行，所以
\begin{align*}
\iint_S xyz \mathrm{d}x\mathrm{d}y &= \iint_{S_1} xyz \mathrm{d}x\mathrm{d}y + \iint_{S_2} xyz \mathrm{d}x\mathrm{d}y \\
&= \iint_{D_{xy}} xy\sqrt{1 - x^2 - y^2} \mathrm{d}x\mathrm{d}y - \iint_{D_{xy}} (-xy\sqrt{1 - x^2 - y^2}) \mathrm{d}x\mathrm{d}y \\
&= 2 \iint_{D_{xy}} xy\sqrt{1 - x^2 - y^2} \mathrm{d}x\mathrm{d}y \\
&= 2 \int_0^{\frac{\pi}{2}} \mathrm{d}\theta \int_0^1 r^3 \cos\theta \sin\theta \sqrt{1 - r^2} \mathrm{d}r = \frac{2}{15}.
\end{align*}
\end{solution}

\begin{proposition}
如果光滑曲面 $S$ 由参数方程给出:
\begin{align*}
S: \begin{cases} x = x(u,v), \\ y = y(u,v), \\ z = z(u,v), \end{cases} \quad (u,v) \in D.
\end{align*}
若在 $D$ 上各点它们的函数行列式
\begin{align*}
\frac{\partial(y,z)}{\partial(u,v)}, \quad \frac{\partial(z,x)}{\partial(u,v)}, \quad \frac{\partial(x,y)}{\partial(u,v)}
\end{align*}
不同时为零，则分别有
\begin{align}
\iint_S P \mathrm{d}y\mathrm{d}z &= \pm \iint_D P(x(u,v), y(u,v), z(u,v)) \frac{\partial(y,z)}{\partial(u,v)} \mathrm{d}u\mathrm{d}v, \label{eq:22-5} \\
\iint_S Q \mathrm{d}z\mathrm{d}x &= \pm \iint_D Q(x(u,v), y(u,v), z(u,v)) \frac{\partial(z,x)}{\partial(u,v)} \mathrm{d}u\mathrm{d}v, \label{eq:22-6} \\
\iint_S R \mathrm{d}x\mathrm{d}y &= \pm \iint_D R(x(u,v), y(u,v), z(u,v)) \frac{\partial(x,y)}{\partial(u,v)} \mathrm{d}u\mathrm{d}v. \label{eq:22-7}
\end{align}
\end{proposition}
\begin{remark}
\eqref{eq:22-5}, \eqref{eq:22-6}, \eqref{eq:22-7} 三式中的正负号分别对应曲面 $S$ 的两个侧，特别当法向 $\left(\frac{\partial(y,z)}{\partial(u,v)}, \frac{\partial(z,x)}{\partial(u,v)}, \frac{\partial(x,y)}{\partial(u,v)}\right)$ 对应于曲面 $S$ 所选定的正向一侧时，取正号，否则取负号.
\end{remark}

\begin{example}
计算 $\iint_S x^3 \mathrm{d}y\mathrm{d}z$，其中 $S$ 为椭球面 $\frac{x^2}{a^2} + \frac{y^2}{b^2} + \frac{z^2}{c^2} = 1$ 的上半部并选取外侧.
\end{example}

\begin{solution}
把曲面表示为参数方程
\begin{align*}
x = a\sin\varphi\cos\theta, \quad y = b\sin\varphi\sin\theta, \quad z = c\cos\varphi \quad \left(0 \leqslant \varphi \leqslant \frac{\pi}{2}, 0 \leqslant \theta \leqslant 2\pi\right).
\end{align*}
由 \eqref{eq:22-5} 式有
\begin{align}
\iint_S x^3 \mathrm{d}y\mathrm{d}z = \pm \iint_{D_{\varphi\theta}} a^3 \sin^3\varphi \cos^3\theta \cdot A \mathrm{d}\varphi\mathrm{d}\theta, \label{eq:22-8}
\end{align}
其中
\begin{align*}
A = \frac{\partial(y,z)}{\partial(\varphi,\theta)} = \begin{vmatrix} b\cos\varphi\sin\theta & b\sin\varphi\cos\theta \\ -c\sin\varphi & 0 \end{vmatrix} = bc\sin^2\varphi\cos\theta,
\end{align*}
积分是在 $S$ 的上侧进行. 由上述的注，\eqref{eq:22-8} 式右端取正号(因为此时法向中 $\frac{\partial(x,y)}{\partial(\varphi,\theta)} = ab\sin\varphi\cos\varphi \geqslant 0$ 对应 $S$ 的正侧)，即
\begin{align*}
\iint_S x^3 \mathrm{d}y\mathrm{d}z &= \iint_{D_{\varphi\theta}} a^3 \sin^3\varphi \cos^3\theta \cdot bc\sin^2\varphi \cos\theta \mathrm{d}\varphi\mathrm{d}\theta \\
&= a^3 bc \int_0^{\frac{\pi}{2}} \sin^5\varphi \mathrm{d}\varphi \int_0^{2\pi} \cos^4\theta \mathrm{d}\theta = \frac{2}{5}\pi a^3 bc.
\end{align*}
\end{solution}

\begin{theorem}
设 $S$ 为光滑曲面，正侧法向量为 $(\cos\alpha, \cos\beta, \cos\gamma)$，$P(x,y,z), Q(x,y,z), R(x,y,z)$ 在 $S$ 上连续，则
\begin{align*}
\iint_S P(x,y,z) \mathrm{d}y\mathrm{d}z + Q(x,y,z) \mathrm{d}z\mathrm{d}x + R(x,y,z) \mathrm{d}x\mathrm{d}y = \iint_S (P(x,y,z)\cos\alpha + Q(x,y,z)\cos\beta + R(x,y,z)\cos\gamma) \mathrm{d}S.
\end{align*}
\end{theorem}

\begin{theorem}
设 $P, Q, R$ 是定义在光滑曲面 $S: z = z(x,y), (x,y) \in D$ 上的连续函数，以 $S$ 的上侧为正侧，则
\begin{align*}
&\iint_S P(x,y,z) \mathrm{d}y\mathrm{d}z + Q(x,y,z) \mathrm{d}z\mathrm{d}x + R(x,y,z) \mathrm{d}x\mathrm{d}y \\
&= \iint_D (P(x,y,z(x,y))(-z_x) + Q(x,y,z(x,y))(-z_y) + R(x,y,z(x,y))) \mathrm{d}x\mathrm{d}y.
\end{align*}
\end{theorem}
\begin{proof}
由于
\begin{align*}
\cos\alpha = \frac{-z_x}{\sqrt{1+z_x^2+z_y^2}}, \quad \cos\beta = \frac{-z_y}{\sqrt{1+z_x^2+z_y^2}}, \quad \cos\gamma = \frac{1}{\sqrt{1+z_x^2+z_y^2}}, \quad \mathrm{d}S = \sqrt{1+z_x^2+z_y^2} \mathrm{d}x\mathrm{d}y,
\end{align*}
因此
\begin{align*}
&\iint_S P(x,y,z) \mathrm{d}y\mathrm{d}z + Q(x,y,z) \mathrm{d}z\mathrm{d}x + R(x,y,z) \mathrm{d}x\mathrm{d}y \\
&= \iint_S (P(x,y,z)\cos\alpha + Q(x,y,z)\cos\beta + R(x,y,z)\cos\gamma) \mathrm{d}S \\
&= \iint_D (P(x,y,z(x,y))(-z_x) + Q(x,y,z(x,y))(-z_y) + R(x,y,z(x,y))) \mathrm{d}x\mathrm{d}y.
\end{align*}
\end{proof}

\begin{example}
计算 $\iint_S (2x+z) \mathrm{d}y\mathrm{d}z + z \mathrm{d}x\mathrm{d}y$，其中 $S = \{(x,y,z) \mid z = x^2 + y^2, z \in [0,1]\}$，取上侧.
\end{example}

\begin{solution}
$z_x = 2x, z_y = 2y$,
\begin{align*}
I &= \iint_D ((2x + x^2 + y^2)(-2x) + x^2 + y^2) \mathrm{d}x\mathrm{d}y = \iint_D (y^2 - 3x^2) \mathrm{d}x\mathrm{d}y \\
&= -\iint_D 2x^2 \mathrm{d}x\mathrm{d}y = -2 \int_0^{2\pi} \mathrm{d}\theta \int_0^1 r^2 \cos^2\theta r \mathrm{d}r = -\frac{\pi}{2},
\end{align*}
其中由于 $x(x^2 + y^2)$ 是 $x$ 的奇函数，$\iint_D x(x^2 + y^2) \mathrm{d}x\mathrm{d}y = 0$；又由对称性，
\begin{align*}
\iint_D x^2 \mathrm{d}x\mathrm{d}y = \iint_D y^2 \mathrm{d}x\mathrm{d}y.
\end{align*}
\end{solution}



\end{document}